\documentclass[12pt,a4paper]{report}

\usepackage[T1]{fontenc}
\usepackage[utf8]{inputenc}
\usepackage[french]{babel}
\usepackage{lmodern}
\usepackage[a4paper,margin=2.5cm]{geometry}
\usepackage{amsmath,amssymb,amsfonts}
\usepackage{siunitx}
\usepackage{graphicx}
\usepackage{booktabs}
\usepackage{array}
\usepackage{float}
\usepackage{hyperref}
\usepackage{enumitem}
\usepackage{xcolor}
\usepackage{listings}
\usepackage{caption}

\hypersetup{
    colorlinks=true,
    linkcolor=blue!50!black,
    urlcolor=blue!50!black,
    citecolor=blue!50!black
}

\lstdefinestyle{pythonstyle}{
    language=Python,
    basicstyle=\ttfamily\small,
    keywordstyle=\color{blue!60!black},
    commentstyle=\color{green!40!black},
    stringstyle=\color{red!50!black},
    showstringspaces=false,
    breaklines=true,
    frame=single
}

\lstdefinestyle{codestyle}{
    basicstyle=\ttfamily\small,
    showstringspaces=false,
    breaklines=true,
    frame=single
}

\newcommand{\Lone}{L_1}
\newcommand{\Ltwo}{L_2}
\newcommand{\Lthree}{L_3}
\newcommand{\thone}{\theta_1}
\newcommand{\thtwo}{\theta_2}
\newcommand{\ththree}{\theta_3}
\newcommand{\ddof}{2~DDL}
\newcommand{\workspace}{espace de travail}
\newcommand{\placeholder}[2]{
    \begin{figure}[H]
        \centering
        \fbox{
            \parbox[c][6cm][c]{0.85\linewidth}{
                \centering
                \textbf{Emplacement réservé pour une capture d'écran}\\[0.5em]
                #2
            }
        }
        \caption{#1}
    \end{figure}
}

\title{
    \Huge Étude, modélisation et implémentation\\[0.4em]
    d'un bras robotique plan à deux degrés de liberté\\[0.4em]
    \Large avec simulation, trajectoires et perspective d'intégration GRBL
}
\author{Projet de thèse}
\date{\today}

\begin{document}

\maketitle
\tableofcontents
\clearpage

\chapter*{Résumé}
\addcontentsline{toc}{chapter}{Résumé}

Ce document présente de manière détaillée le projet de réalisation d'un bras robotique
plan à deux degrés de liberté, commandé en position à partir d'un modèle cinématique et
destiné à évoluer vers une validation expérimentale sur prototype physique. Le travail
réunit trois dimensions complémentaires : la théorie mathématique, l'implémentation
logicielle et la préparation à l'intégration matérielle.

Le coeur du projet repose sur la cinématique directe et inverse d'un manipulateur plan
de type $RR$, c'est-à-dire deux articulations rotoïdes en série. Le modèle obtenu permet
de passer d'une représentation articulaire $(\thone,\thtwo)$ à une représentation
cartésienne $(x,y)$, puis d'inverser cette relation pour atteindre un point cible.
Sur cette base, le projet met en oeuvre plusieurs modes de fonctionnement : une
simulation graphique, une exécution de trajectoires et un mode hybride combinant
visualisation et dialogue avec une chaîne de commande inspirée de GRBL.

L'objectif final n'est pas seulement de déplacer un bras dans une animation, mais de
montrer la cohérence entre les modèles mathématiques, les décisions logicielles et la
future validation matérielle. Cette cohérence est essentielle dans un contexte de thèse,
car elle permet de démontrer que l'implémentation numérique constitue déjà une preuve de
faisabilité avant la fabrication complète du prototype.

\chapter{Introduction générale}

\section{Contexte}

Les bras robotiques à faible nombre de degrés de liberté constituent un excellent cadre
d'étude pour aborder la robotique de manière rigoureuse. Ils permettent d'introduire
les notions fondamentales de cinématique, d'espace de travail, de singularité, de
génération de trajectoires et de commande, tout en conservant une complexité compatible
avec une première validation expérimentale.

Dans ce projet, le système étudié est un bras robotique plan à deux degrés de liberté.
Le bras comprend deux segments rigides, articulés autour de deux axes de rotation. Son
effecteur final évolue dans le plan $(x,y)$, ce qui permet de relier naturellement les
variables articulaires aux coordonnées cartésiennes d'un point à atteindre.

\section{Problématique}

La problématique générale peut être formulée ainsi :

\begin{quote}
Comment concevoir un système complet permettant de calculer les angles articulaires
nécessaires pour atteindre une position donnée, de visualiser de manière crédible ce
mouvement, puis de préparer le passage à une exécution réelle sur prototype commandé
par GRBL ?
\end{quote}

Cette problématique se décompose en quatre sous-questions :

\begin{enumerate}[label=\arabic*.]
    \item Comment établir un modèle mathématique fiable du bras robotique ?
    \item Comment transformer ce modèle en logiciel exécutable et observable ?
    \item Comment structurer l'architecture pour préparer une liaison avec du matériel réel ?
    \item Comment faire évoluer le système vers un troisième degré de liberté ?
\end{enumerate}

\section{Objectifs du projet}

Les objectifs poursuivis sont les suivants :

\begin{itemize}
    \item établir la cinématique directe et inverse d'un bras plan \ddof ;
    \item modéliser l'\workspace{} et les singularités du robot ;
    \item implémenter une simulation graphique animée du mouvement ;
    \item générer et exécuter des trajectoires cartésiennes ;
    \item préparer une architecture logicielle compatible avec une chaîne matérielle pilotée par GRBL ;
    \item définir une perspective claire d'extension vers un robot à trois degrés de liberté.
\end{itemize}

\section{Organisation du document}

Le présent document suit une progression naturelle :

\begin{itemize}
    \item le chapitre \ref{chap:systeme} présente le système étudié ;
    \item le chapitre \ref{chap:theorie} développe les modèles mathématiques ;
    \item le chapitre \ref{chap:commande} décrit la génération des mouvements ;
    \item le chapitre \ref{chap:implementation} détaille l'implémentation logicielle ;
    \item le chapitre \ref{chap:grbl} introduit l'intégration future de GRBL sur prototype ;
    \item le chapitre \ref{chap:perspectives} expose l'extension vers un troisième degré de liberté.
\end{itemize}

\chapter{Présentation du système}\label{chap:systeme}

\section{Description mécanique}

Le robot considéré est un manipulateur plan à deux articulations rotoïdes. Il est formé :

\begin{itemize}
    \item d'une base fixe ;
    \item d'un premier segment de longueur $\Lone$ ;
    \item d'un second segment de longueur $\Ltwo$ ;
    \item d'un effecteur final situé à l'extrémité du second segment.
\end{itemize}

Dans l'état actuel du projet, les paramètres utilisés dans le logiciel sont :

\begin{table}[H]
    \centering
    \caption{Paramètres principaux du bras robotique}
    \begin{tabular}{>{\raggedright\arraybackslash}p{7cm}cc}
        \toprule
        Paramètre & Symbole & Valeur \\
        \midrule
        Longueur du premier segment & $\Lone$ & \SI{200}{mm} \\
        Longueur du second segment & $\Ltwo$ & \SI{150}{mm} \\
        Limite articulaire 1 & $\thone$ & $[-\pi,\pi]$ \\
        Limite articulaire 2 & $\thtwo$ & $\left[-150^{\circ},150^{\circ}\right]$ \\
        Pas moteur par révolution & -- & $200$ \\
        Micro-pas & -- & $16$ \\
        Fréquence d'animation & -- & \SI{60}{fps} \\
        \bottomrule
    \end{tabular}
\end{table}

\section{Repères et conventions}

Le repère de travail est choisi de la manière suivante :

\begin{itemize}
    \item l'origine est placée à la base du robot ;
    \item l'axe $x$ est horizontal et orienté vers la droite ;
    \item l'axe $y$ est vertical et orienté vers le haut ;
    \item les angles sont comptés dans le sens trigonométrique.
\end{itemize}

L'articulation 1 correspond à la rotation de base, tandis que l'articulation 2
correspond à la rotation du coude.

\placeholder
{Schéma cinématique du bras robotique \ddof}
{Insérer ici un schéma ou une capture représentant la base, les deux segments,
les angles $\theta_1$ et $\theta_2$, ainsi que l'effecteur final.}

\section{Hypothèses de modélisation}

Pour conserver un cadre analytique clair, les hypothèses suivantes sont adoptées :

\begin{itemize}
    \item les segments sont supposés rigides ;
    \item les liaisons sont parfaites et sans jeu dans la modélisation théorique ;
    \item le mouvement est purement plan ;
    \item les effets dynamiques, les frottements et les couples ne sont pas encore intégrés dans la loi de commande principale ;
    \item le contrôle est actuellement de type position, en boucle ouverte, avec validation d'abord par simulation.
\end{itemize}

Ces hypothèses sont adaptées à une première phase de validation, où l'objectif
principal est la cohérence cinématique et la faisabilité expérimentale.

\chapter{Fondements théoriques et modèles mathématiques}\label{chap:theorie}

\section{Cinématique directe}

\subsection{Principe}

La cinématique directe consiste à calculer la position cartésienne de l'effecteur final
à partir des angles articulaires. Pour un bras plan \ddof, cette relation est obtenue
par composition vectorielle des deux segments.

\subsection{Équations}

La position du premier joint intermédiaire est :

\begin{align}
    x_1 &= \Lone \cos(\thone), \\
    y_1 &= \Lone \sin(\thone).
\end{align}

La position de l'effecteur final devient alors :

\begin{align}
    x &= \Lone \cos(\thone) + \Ltwo \cos(\thone + \thtwo), \label{eq:fkx} \\
    y &= \Lone \sin(\thone) + \Ltwo \sin(\thone + \thtwo). \label{eq:fky}
\end{align}

Ces expressions constituent le modèle central de la cinématique directe.

\subsection{Interprétation}

La première contribution correspond à la projection du premier bras sur les axes du
repère. La seconde contribution correspond à la projection du second bras, dont
l'orientation absolue est la somme des deux angles articulaires.

\subsection{Implémentation}

Dans l'application, cette partie est prise en charge par le module
\texttt{phase1\_simulation/kinematics.py}, au travers de la méthode
\texttt{forward\_kinematics}. Ce calcul est utilisé :

\begin{itemize}
    \item pour afficher la position courante ;
    \item pour dessiner le bras dans le simulateur ;
    \item pour tracer l'historique de trajectoire ;
    \item pour vérifier la cohérence de la cinématique inverse.
\end{itemize}

\section{Cinématique inverse}

\subsection{Principe}

La cinématique inverse consiste à déterminer les angles $(\thone,\thtwo)$ permettant
d'atteindre une cible cartésienne $(x,y)$. Cette opération est plus délicate que la
cinématique directe, car elle peut conduire à plusieurs solutions ou à aucune solution.

\subsection{Calcul de la distance à la cible}

La distance entre la base et le point cible vaut :

\begin{equation}
    r = \sqrt{x^2 + y^2}.
\end{equation}

Cette grandeur permet de savoir immédiatement si la cible appartient à l'\workspace{}.

\subsection{Condition d'atteignabilité}

Le point est atteignable si et seulement si :

\begin{equation}
    |\Lone - \Ltwo| \leq r \leq \Lone + \Ltwo.
\end{equation}

Pour les valeurs actuelles du projet, on obtient :

\begin{align}
    R_{\max} &= \Lone + \Ltwo = \SI{350}{mm}, \\
    R_{\min} &= |\Lone - \Ltwo| = \SI{50}{mm}.
\end{align}

\subsection{Calcul de l'angle du coude}

L'application de la loi des cosinus fournit :

\begin{equation}
    \cos(\thtwo) = \frac{r^2 - \Lone^2 - \Ltwo^2}{2 \Lone \Ltwo}.
\end{equation}

On en déduit :

\begin{equation}
    \thtwo = \pm \arccos\left(\frac{r^2 - \Lone^2 - \Ltwo^2}{2 \Lone \Ltwo}\right).
\end{equation}

Les deux signes correspondent aux deux configurations classiques :

\begin{itemize}
    \item coude vers le haut ;
    \item coude vers le bas.
\end{itemize}

\subsection{Calcul de l'angle de base}

Une fois $\thtwo$ déterminé, l'angle de base s'écrit :

\begin{equation}
    \thone = \operatorname{atan2}(y,x)
    - \operatorname{atan2}\left(\Ltwo \sin(\thtwo), \Lone + \Ltwo \cos(\thtwo)\right).
\end{equation}

\subsection{Exemple numérique}

Considérons la cible $(x,y) = (250,150)$~mm avec $\Lone=\SI{200}{mm}$ et
$\Ltwo=\SI{150}{mm}$.

\begin{align}
    r &= \sqrt{250^2 + 150^2} = \SI{291.55}{mm}, \\
    \cos(\thtwo) &= \frac{291.55^2 - 200^2 - 150^2}{2 \times 200 \times 150}
    \approx 0.374.
\end{align}

On en déduit une solution de type coude vers le haut de l'ordre de :

\begin{align}
    \thtwo &\approx 68.0^\circ, \\
    \thone &\approx 2.5^\circ.
\end{align}

Cette solution est cohérente avec l'implémentation actuellement validée dans le code.

\subsection{Implémentation logicielle}

L'algorithme de cinématique inverse est encapsulé dans la méthode
\texttt{inverse\_kinematics}. Le logiciel :

\begin{itemize}
    \item vérifie d'abord l'atteignabilité du point ;
    \item calcule ensuite $\cos(\thtwo)$ ;
    \item borne numériquement cette valeur dans l'intervalle $[-1,1]$ ;
    \item choisit l'une des deux branches ;
    \item vérifie enfin les limites articulaires.
\end{itemize}

\section{Espace de travail}

L'\workspace{} du bras est une couronne circulaire comprise entre les rayons
$R_{\min}$ et $R_{\max}$.

Cette géométrie a une conséquence importante pour l'interface utilisateur :
un point visuellement situé trop près de la base ne doit pas être considéré comme
atteignable, tout comme un point situé à l'extérieur du rayon maximum.

L'affichage graphique du projet représente explicitement cette couronne afin que
l'utilisateur puisse interpréter visuellement les zones accessibles et inaccessibles.

\placeholder
{Visualisation de l'espace de travail}
{Insérer ici une capture du simulateur montrant le cercle externe, le cercle interne
et la position du bras dans la zone accessible.}

\section{Jacobienne}

\subsection{Définition}

La matrice jacobienne relie les vitesses articulaires aux vitesses cartésiennes :

\begin{equation}
    \begin{bmatrix}
        \dot{x} \\
        \dot{y}
    \end{bmatrix}
    =
    J(\thone,\thtwo)
    \begin{bmatrix}
        \dot{\theta}_1 \\
        \dot{\theta}_2
    \end{bmatrix}.
\end{equation}

\subsection{Expression analytique}

En dérivant les équations \eqref{eq:fkx} et \eqref{eq:fky}, on obtient :

\begin{equation}
    J(\thone,\thtwo) =
    \begin{bmatrix}
        -\Lone \sin(\thone) - \Ltwo \sin(\thone+\thtwo)
        &
        -\Ltwo \sin(\thone+\thtwo)
        \\
        \Lone \cos(\thone) + \Ltwo \cos(\thone+\thtwo)
        &
        \Ltwo \cos(\thone+\thtwo)
    \end{bmatrix}.
\end{equation}

\subsection{Rôle dans le projet}

Même si le projet n'utilise pas encore la jacobienne comme outil principal de commande,
sa présence est importante pour trois raisons :

\begin{itemize}
    \item analyser les singularités ;
    \item préparer une évolution vers une commande en vitesse ou en effort ;
    \item fournir un socle théorique solide pour la suite de la thèse.
\end{itemize}

\section{Singularités}

\subsection{Détection}

Une singularité apparaît lorsque le déterminant de la jacobienne est nul :

\begin{equation}
    \det(J) = \Lone \Ltwo \sin(\thtwo) = 0.
\end{equation}

Donc les configurations singulières apparaissent lorsque :

\begin{equation}
    \thtwo = 0 \quad \text{ou} \quad \thtwo = \pi.
\end{equation}

\subsection{Signification physique}

Ces cas correspondent à des configurations alignées, soit complètement déployées,
soit complètement repliées. Dans ces situations, certaines directions de mouvement
cartésien deviennent difficiles, voire impossibles, à produire localement.

\subsection{Conséquence pratique}

La prise en compte des singularités est importante pour la planification de trajectoire,
car une trajectoire mathématiquement valide peut devenir défavorable en pratique si elle
s'approche trop d'une zone où les vitesses articulaires devraient devenir très élevées.

\chapter{Génération des mouvements et stratégie de commande}\label{chap:commande}

\section{Commande d'un point cible}

Le cas le plus simple consiste à atteindre un point cible $(x_c,y_c)$. La chaîne de
traitement est alors :

\begin{enumerate}[label=\arabic*.]
    \item saisie ou sélection de la cible ;
    \item calcul de la cinématique inverse ;
    \item validation des limites ;
    \item mise à jour des consignes articulaires ;
    \item interpolation du mouvement pour l'animation ;
    \item exécution simulée, matérielle ou hybride.
\end{enumerate}

\section{Trajectoire linéaire dans l'espace cartésien}

Pour relier un point initial $(x_0,y_0)$ à un point final $(x_f,y_f)$, le projet
met en oeuvre une interpolation cartésienne discrète :

\begin{align}
    x_i &= x_0 + \frac{i}{N-1}(x_f-x_0), \\
    y_i &= y_0 + \frac{i}{N-1}(y_f-y_0),
\end{align}

où $N$ est le nombre de points discrétisés.

Pour chaque point intermédiaire, la cinématique inverse est recalculée. Cette stratégie
présente deux intérêts :

\begin{itemize}
    \item le chemin suivi dans le plan est conforme à l'intuition visuelle ;
    \item la simulation devient plus démonstrative qu'un saut direct vers la position finale.
\end{itemize}

\section{Trajectoires de démonstration}

Au-delà du déplacement vers un point, plusieurs trajectoires de démonstration ont été
structurées dans le projet :

\begin{itemize}
    \item trajectoire linéaire ;
    \item trajectoire carrée ;
    \item trajectoire circulaire ;
    \item trajectoire en zigzag.
\end{itemize}

Ces trajectoires servent un double objectif :

\begin{itemize}
    \item valider la robustesse du calcul cinématique sur plusieurs points successifs ;
    \item fournir des résultats visuels convaincants pour une démonstration ou une soutenance.
\end{itemize}

\section{Interpolation articulaire pour l'animation}

Pour qu'un mouvement soit observable et crédible, le projet n'effectue pas une simple
mise à jour instantanée des angles. Une interpolation temporelle est introduite entre
l'état initial et l'état final.

Si $\theta_k^0$ et $\theta_k^f$ désignent les angles initial et final d'une articulation
$k \in \{1,2\}$, alors l'évolution peut être écrite sous la forme :

\begin{equation}
    \theta_k(t) = \theta_k^0 + s(t)\left(\theta_k^f - \theta_k^0\right),
\end{equation}

où $s(t)$ est une fonction de lissage vérifiant $s(0)=0$ et $s(T)=1$.

Dans l'implémentation actuelle, une loi de type cosinus est utilisée :

\begin{equation}
    s(\alpha) = \frac{1}{2} - \frac{1}{2}\cos(\pi \alpha),
    \qquad \alpha \in [0,1].
\end{equation}

Cette écriture évite un déplacement visuellement brutal au démarrage et à l'arrivée.

\section{Nature de la commande}

Le projet relève d'une commande en position, principalement en boucle ouverte.
Autrement dit, la consigne est calculée à partir du modèle, puis transmise au système
de simulation ou au futur système matériel sans mesure continue de correction dans la
boucle principale.

Ce choix est adapté à une phase de prototypage académique, car il permet :

\begin{itemize}
    \item de valider la chaîne de calcul ;
    \item d'identifier les erreurs de modélisation ;
    \item de préparer une future extension vers une commande plus élaborée si nécessaire.
\end{itemize}

\chapter{Implémentation logicielle du projet}\label{chap:implementation}

\section{Vue d'ensemble}

Le dépôt est organisé en trois niveaux fonctionnels :

\begin{enumerate}[label=\arabic*.]
    \item simulation et modélisation ;
    \item commande matérielle ;
    \item intégration complète.
\end{enumerate}

\begin{lstlisting}[style=codestyle,caption={Organisation générale du projet}]
phase1_simulation/
    config.py
    kinematics.py
    simulator.py

phase2_hardware/
    grbl_interface.py
    motor_control.py
    a4988_config.py

phase3_integration/
    robot_controller.py
    gui.py

tests/
    test_kinematics.py
    test_robot_controller.py
\end{lstlisting}

\section{Phase 1 : configuration et modélisation}

\subsection{Module \texttt{config.py}}

Ce module centralise les paramètres géométriques, angulaires, moteurs et graphiques.
Il sert de point unique de cohérence pour l'ensemble du projet. Cette approche évite
la dispersion des constantes dans le code et facilite les recalibrages.

\subsection{Module \texttt{kinematics.py}}

Ce module implémente :

\begin{itemize}
    \item la cinématique directe ;
    \item la cinématique inverse ;
    \item la jacobienne ;
    \item le calcul de la frontière de l'espace de travail ;
    \item la génération de trajectoires linéaires ;
    \item la génération de trajectoires polylignes et circulaires.
\end{itemize}

Il constitue donc le noyau mathématique du projet.

\subsection{Module \texttt{simulator.py}}

Le simulateur graphique en \texttt{Pygame} représente :

\begin{itemize}
    \item la base du robot ;
    \item les deux segments ;
    \item l'effecteur final ;
    \item la cible ;
    \item l'historique de trajectoire ;
    \item l'\workspace{} sous forme de couronne.
\end{itemize}

L'évolution récente du projet a consisté à rendre cette simulation véritablement animée.
Le bras n'est plus figé ni téléporté : les angles sont interpolés sur une durée finie,
ce qui améliore la lisibilité et la valeur démonstrative du système.

\placeholder
{Simulation graphique du mouvement vers un point cible}
{Insérer ici une capture montrant l'effecteur, la cible et l'historique du mouvement
pendant une animation d'atteinte de point.}

\section{Phase 2 : préparation matérielle}

\subsection{Module \texttt{grbl\_interface.py}}

Ce module encapsule la communication série avec GRBL. Il gère notamment :

\begin{itemize}
    \item l'ouverture et la fermeture du port série ;
    \item l'envoi de commandes G-code ;
    \item la lecture des réponses de type \texttt{ok} ou \texttt{error:N} ;
    \item la lecture du statut de la carte ;
    \item les commandes de homing, de reset et de déverrouillage.
\end{itemize}

\subsection{Module \texttt{motor\_control.py}}

Ce module constitue une couche d'abstraction entre les angles articulaires et les
commandes envoyées à GRBL. Son rôle est important, car il matérialise la transition
entre le monde mathématique et le monde électromécanique.

\subsection{Module \texttt{a4988\_config.py}}

Ce module regroupe les éléments de configuration liés au driver A4988 :

\begin{itemize}
    \item modes de micro-pas ;
    \item réglage du courant ;
    \item rappels de câblage ;
    \item paramètres utiles pour la préparation du prototype.
\end{itemize}

\section{Phase 3 : intégration complète}

\subsection{Module \texttt{robot\_controller.py}}

Le contrôleur principal fédère les trois modes du projet :

\begin{itemize}
    \item \textbf{simulation} : calcul et animation sans matériel ;
    \item \textbf{hardware} : envoi des commandes vers GRBL ;
    \item \textbf{hybrid} : animation et exécution matérielle simultanées.
\end{itemize}

Ce module gère également :

\begin{itemize}
    \item l'état courant du robot ;
    \item l'historique des positions ;
    \item les trajectoires planifiées ;
    \item les rappels d'état destinés à l'interface graphique ;
    \item l'animation temporelle des mouvements.
\end{itemize}

\subsection{Module \texttt{gui.py}}

L'interface \texttt{Tkinter} permet de piloter le robot à un niveau plus démonstratif.
Elle intègre :

\begin{itemize}
    \item un contrôle par angles ;
    \item un contrôle par position cartésienne ;
    \item une planification de trajectoire ;
    \item des trajectoires prédéfinies ;
    \item une console de messages ;
    \item une vue temps réel du bras.
\end{itemize}

L'intérêt scientifique de cette interface est majeur : elle ne sert pas seulement
à piloter, mais aussi à rendre visible la correspondance entre la théorie, le code
et le comportement du robot.

\placeholder
{Interface graphique intégrée du projet}
{Insérer ici une capture de l'interface complète, en particulier la vue temps réel
du bras et les panneaux de commande.}

\section{Flux de données}

Le flux de données peut être résumé ainsi :

\begin{enumerate}[label=\arabic*.]
    \item l'utilisateur demande un mouvement ;
    \item le contrôleur interprète la demande ;
    \item la cinématique inverse calcule les angles requis ;
    \item le contrôleur anime l'évolution articulaire ;
    \item en mode matériel ou hybride, une couche convertit la consigne en dialogue GRBL ;
    \item l'interface restitue l'état courant et la trajectoire.
\end{enumerate}

\section{Validation logicielle}

Le projet contient des tests unitaires dédiés aux modules critiques, notamment :

\begin{itemize}
    \item tests de cinématique directe ;
    \item tests de cinématique inverse ;
    \item tests de cohérence angles/pas ;
    \item tests sur la planification et l'exécution de trajectoires.
\end{itemize}

Ces tests jouent un rôle fondamental : ils garantissent que la démonstration visuelle
repose sur un comportement mathématique vérifié, et non sur une simple animation.

\chapter{Introduction à GRBL dans la phase prototype}\label{chap:grbl}

\section{Pourquoi GRBL ?}

GRBL est un micrologiciel libre largement utilisé pour piloter des moteurs pas à pas
à partir de commandes G-code. Il est historiquement conçu pour les machines à commande
numérique, mais son principe de fonctionnement est suffisamment générique pour être
adapté à un bras robotique de petite taille, à condition de définir une couche
d'interprétation appropriée.

Dans ce projet, GRBL est envisagé comme une solution pragmatique pour le prototype,
car il fournit déjà :

\begin{itemize}
    \item la communication série ;
    \item la gestion des vitesses d'avance ;
    \item la génération de pas ;
    \item les mécanismes de homing et de sécurité ;
    \item un protocole robuste et documenté.
\end{itemize}

\section{Principe d'intégration}

L'idée n'est pas de faire croire à GRBL qu'il pilote directement un bras dans l'espace
cartésien du plan, mais de l'utiliser comme couche de bas niveau pour produire les
déplacements des moteurs correspondant aux angles articulaires.

Le principe retenu est le suivant :

\begin{enumerate}[label=\arabic*.]
    \item le modèle cinématique calcule $(\thone,\thtwo)$ ;
    \item ces angles sont convertis en nombres de pas moteurs ;
    \item une relation de calibration convertit ces pas en positions compatibles avec les axes GRBL ;
    \item GRBL reçoit alors une commande sur ses axes \texttt{X} et \texttt{Y}.
\end{enumerate}

\section{Conversion angles vers pas}

Soit $N_1$ et $N_2$ les nombres de pas équivalents aux angles articulaires :

\begin{align}
    N_1 &= \thone \, S_1, \\
    N_2 &= \thtwo \, S_2,
\end{align}

où :

\begin{align}
    S_1 &= \frac{N_{\text{rev}} \, \mu_1 \, g_1}{2\pi}, \\
    S_2 &= \frac{N_{\text{rev}} \, \mu_2 \, g_2}{2\pi}.
\end{align}

Ici :

\begin{itemize}
    \item $N_{\text{rev}}$ est le nombre de pas par révolution du moteur ;
    \item $\mu_i$ est le facteur de micro-pas ;
    \item $g_i$ est le rapport de réduction mécanique éventuel.
\end{itemize}

Dans l'état logiciel actuel, cette conversion est déjà présente via
\texttt{steps\_per\_rad\_1} et \texttt{steps\_per\_rad\_2}.

\section{Abstraction vers les axes GRBL}

Une couche de mapping associe :

\begin{itemize}
    \item l'axe \texttt{X} de GRBL à l'articulation 1 ;
    \item l'axe \texttt{Y} de GRBL à l'articulation 2.
\end{itemize}

Dans une première version, on peut poser :

\begin{align}
    X_{\text{grbl}} &= k_s N_1, \\
    Y_{\text{grbl}} &= k_s N_2,
\end{align}

où $k_s$ est un coefficient de calibration exprimé en millimètres GRBL par pas moteur.

Il est important de noter que ce choix relève d'une abstraction d'intégration et non
d'une vérité géométrique universelle. Lors du prototype physique, cette relation devra
être calibrée expérimentalement afin que la consigne envoyée corresponde exactement au
déplacement angulaire réel.

\section{Exemples de commandes G-code}

Une commande de déplacement absolu peut prendre la forme :

\begin{lstlisting}[style=codestyle,caption={Exemple de commande G-code pour une consigne articulaire}]
G90 G1 X12.500 Y8.200 F1000
\end{lstlisting}

où :

\begin{itemize}
    \item \texttt{G90} active le mode absolu ;
    \item \texttt{G1} demande une interpolation linéaire ;
    \item \texttt{X} et \texttt{Y} représentent ici les deux articulations abstraites ;
    \item \texttt{F} impose la vitesse d'avance.
\end{itemize}

D'autres commandes importantes sont :

\begin{lstlisting}[style=codestyle,caption={Commandes utiles pour la phase prototype}]
$H      % homing
$X      % unlock apres alarme
?       % lecture du statut
\end{lstlisting}

\section{Rôle du homing}

Sur le prototype, le homing sera indispensable pour référencer la machine. Sans cette
opération, aucune cohérence durable entre la simulation et le robot réel ne peut être
garantie.

Le homing permettra :

\begin{itemize}
    \item de retrouver une origine articulaire connue ;
    \item de limiter la dérive cumulative des positions ;
    \item de sécuriser les essais ;
    \item de comparer les expériences dans des conditions reproductibles.
\end{itemize}

\section{Paramètres GRBL à considérer}

Les paramètres les plus sensibles dans la phase prototype seront :

\begin{itemize}
    \item les pas par millimètre abstrait des axes \texttt{\$100} et \texttt{\$101} ;
    \item les vitesses maximales \texttt{\$110} et \texttt{\$111} ;
    \item les accélérations \texttt{\$120} et \texttt{\$121} ;
    \item l'activation du homing ;
    \item les limites logicielles et matérielles.
\end{itemize}

\section{Protocole expérimental recommandé pour le prototype}

Lorsque le prototype sera disponible, une séquence de validation progressive est
recommandée :

\begin{enumerate}[label=\arabic*.]
    \item vérifier le câblage et le réglage des courants des A4988 ;
    \item établir la communication série avec GRBL ;
    \item effectuer un homing fiable ;
    \item commander de faibles mouvements articulaires ;
    \item mesurer la correspondance entre consigne et déplacement réel ;
    \item calibrer les coefficients de conversion ;
    \item tester l'atteinte d'un point ;
    \item tester ensuite une trajectoire complète en mode hybride.
\end{enumerate}

\placeholder
{Prototype physique et intégration GRBL}
{Insérer ici une photographie du prototype, du câblage et éventuellement une capture
de la console de communication GRBL.}

\section{Limites actuelles et intérêt scientifique}

À ce stade, l'intégration GRBL doit être comprise comme une passerelle vers le prototype.
Le système démontre déjà la logique complète : calcul d'angles, animation, préparation
des trajectoires, abstraction vers des commandes matérielles. Cette approche est
scientifiquement pertinente, car elle isole clairement les niveaux de complexité :

\begin{itemize}
    \item la géométrie du problème ;
    \item la logique de commande ;
    \item l'exécution électromécanique.
\end{itemize}

\chapter{Analyse des résultats à présenter}\label{chap:resultats}

\section{Résultats attendus en simulation}

La validation visuelle doit montrer que :

\begin{itemize}
    \item le bras suit réellement la cible au lieu de rester figé ;
    \item la trajectoire affichée correspond au chemin voulu ;
    \item la distinction entre mode point, mode trajectoire et mode hybride est visible ;
    \item l'historique de mouvement permet d'interpréter la cohérence de l'exécution.
\end{itemize}

\section{Captures d'écran recommandées}

Pour documenter efficacement le travail, il sera utile d'insérer :

\begin{enumerate}[label=\arabic*.]
    \item une capture du simulateur au repos avec l'\workspace{} visible ;
    \item une capture d'une atteinte de point ;
    \item une capture d'une trajectoire carrée ou circulaire ;
    \item une capture de l'interface graphique complète ;
    \item une capture du mode hybride lors de l'échange avec GRBL ;
    \item une photographie du prototype lorsqu'il sera disponible.
\end{enumerate}

\section{Lectures possibles des résultats}

Dans l'analyse, il conviendra d'insister sur le fait que la simulation n'est pas une
simple illustration décorative. Elle constitue un outil de validation scientifique,
car elle permet de vérifier :

\begin{itemize}
    \item la cohérence des équations ;
    \item l'impact des limites articulaires ;
    \item la forme de l'\workspace{} ;
    \item le comportement en voisinage de trajectoires plus complexes.
\end{itemize}

\chapter{Perspective d'extension vers un troisième degré de liberté}\label{chap:perspectives}

\section{Motivation}

L'ajout d'un troisième degré de liberté constitue la suite naturelle du projet. Un bras
\ddof{} permet de positionner l'effecteur dans un plan, mais ne contrôle pas son
orientation de manière indépendante. Or, dans un grand nombre d'applications robotiques,
la position seule ne suffit pas : il faut également maîtriser l'orientation de l'outil.

\section{Deux scénarios possibles}

Deux évolutions sont particulièrement naturelles :

\begin{itemize}
    \item l'ajout d'une troisième articulation rotoïde pour contrôler l'orientation terminale ;
    \item l'ajout d'un axe prismatic ou vertical si l'application exige une variation de hauteur.
\end{itemize}

Dans la continuité la plus directe du projet actuel, le scénario le plus cohérent est
l'ajout d'une troisième articulation rotoïde, ce qui conduit à un bras plan de type $RRR$.

\section{Extension du modèle cinématique}

Avec une troisième articulation de longueur $\Lthree$ et d'angle $\ththree$, la position
de l'effecteur devient :

\begin{align}
    x &= \Lone \cos(\thone)
    + \Ltwo \cos(\thone+\thtwo)
    + L_3 \cos(\thone+\thtwo+\ththree), \\
    y &= \Lone \sin(\thone)
    + \Ltwo \sin(\thone+\thtwo)
    + L_3 \sin(\thone+\thtwo+\ththree).
\end{align}

Si l'on souhaite en plus imposer une orientation terminale $\phi$, on obtient :

\begin{equation}
    \phi = \thone + \thtwo + \ththree.
\end{equation}

Le problème n'est alors plus seulement de positionner l'effecteur, mais aussi de fixer
son attitude finale.

\section{Impact sur la cinématique inverse}

Le passage au cas $RRR$ augmente nettement la richesse du problème :

\begin{itemize}
    \item il devient possible d'imposer simultanément position et orientation ;
    \item le nombre de solutions peut augmenter ;
    \item la gestion des limites articulaires devient plus délicate ;
    \item la jacobienne devient une matrice $3 \times 3$ si l'orientation est intégrée dans les sorties.
\end{itemize}

\section{Impact sur l'architecture logicielle}

L'ajout du troisième degré de liberté impliquera de faire évoluer :

\begin{itemize}
    \item \texttt{config.py} : ajout des paramètres de la troisième articulation ;
    \item \texttt{kinematics.py} : nouvelles équations de cinématique directe, inverse et jacobienne ;
    \item \texttt{simulator.py} : dessin d'un troisième segment ;
    \item \texttt{robot\_controller.py} : animation et trajectoires à trois variables articulaires ;
    \item \texttt{gui.py} : nouveaux contrôles et nouvelle lecture d'état ;
    \item la couche GRBL : ajout d'un troisième axe abstrait, typiquement \texttt{Z} ou \texttt{A}.
\end{itemize}

\section{Impact sur le prototype}

Sur le plan matériel, cette évolution demandera :

\begin{itemize}
    \item un moteur supplémentaire ;
    \item un driver supplémentaire ;
    \item une adaptation de la structure mécanique ;
    \item une révision du câblage et de l'alimentation ;
    \item une nouvelle phase de calibration.
\end{itemize}

\section{Intérêt scientifique de cette perspective}

L'intérêt de cette extension est double :

\begin{itemize}
    \item elle transforme le projet en système plus proche d'un manipulateur utile pour la préhension ou l'orientation d'outil ;
    \item elle ouvre la voie à des études plus avancées sur la redondance, les contraintes de trajectoire, la commande et l'optimisation.
\end{itemize}

Cette perspective donne au projet une continuité scientifique forte et justifie
pleinement son inscription dans un travail de recherche.

\chapter{Conclusion générale}

Le projet présenté constitue une base cohérente et solide pour l'étude d'un bras
robotique plan à deux degrés de liberté. Sa valeur ne réside pas seulement dans
l'existence d'une interface ou d'une animation, mais dans l'alignement entre :

\begin{itemize}
    \item une théorie mathématique explicite ;
    \item une implémentation logicielle modulaire ;
    \item une stratégie de visualisation démonstrative ;
    \item une préparation concrète à l'intégration sur prototype via GRBL.
\end{itemize}

Les résultats actuels montrent que le système est capable :

\begin{itemize}
    \item de calculer correctement les angles pour atteindre une cible ;
    \item de représenter de manière animée les mouvements du bras ;
    \item d'exécuter des trajectoires plus riches qu'un simple déplacement ponctuel ;
    \item de préparer une future validation expérimentale cohérente.
\end{itemize}

L'étape suivante, à savoir la construction du prototype puis l'ajout d'un troisième
degré de liberté, n'apparaît donc pas comme une rupture, mais comme la continuité
logique d'une architecture déjà pensée pour évoluer.

\appendix

\chapter{Extrait de code représentatif}

\begin{lstlisting}[style=pythonstyle,caption={Exemple simplifié de cinématique directe}]
def forward_kinematics(theta1, theta2, L1, L2):
    x = L1 * np.cos(theta1) + L2 * np.cos(theta1 + theta2)
    y = L1 * np.sin(theta1) + L2 * np.sin(theta1 + theta2)
    return x, y
\end{lstlisting}

\begin{lstlisting}[style=pythonstyle,caption={Exemple simplifié d'interpolation cartésienne}]
def compute_trajectory(start_pos, end_pos, num_points):
    x_traj = np.linspace(start_pos[0], end_pos[0], num_points)
    y_traj = np.linspace(start_pos[1], end_pos[1], num_points)
    return list(zip(x_traj, y_traj))
\end{lstlisting}

\chapter{Références bibliographiques}

\begin{thebibliography}{9}

\bibitem{spong}
M. W. Spong, S. Hutchinson, and M. Vidyasagar,
\textit{Robot Modeling and Control},
Wiley, 2006.

\bibitem{craig}
J. J. Craig,
\textit{Introduction to Robotics: Mechanics and Control},
Pearson, 3rd edition, 2005.

\bibitem{siciliano}
B. Siciliano, L. Sciavicco, L. Villani, and G. Oriolo,
\textit{Robotics: Modelling, Planning and Control},
Springer, 2009.

\bibitem{grbl}
GRBL Project,
\textit{GRBL Wiki},
\url{https://github.com/gnea/grbl/wiki}.

\bibitem{a4988}
Pololu,
\textit{A4988 Stepper Motor Driver Carrier Documentation},
\url{https://www.pololu.com/product/1182}.

\end{thebibliography}

\end{document}
